% =========================================================================
% ABSTRACT (Sommario)
% =========================================================================

\chapter*{Sommario}
\addcontentsline{toc}{chapter}{Sommario}

Molti dei sistemi che caratterizzano il mondo intorno a noi sono invisibili, oppure così complessi --- dominati da innumerevoli variabili spesso dipendenti le une dalle altre --- da resistere a ogni previsione, o entrambe le cose insieme. Si muovono sotto la superficie, si dispiegano su scale temporali più lunghe di una vita umana e non si lasciano leggere in un singolo istante. La difficoltà non è la mancanza di dati, ma lo scarto tra dati e comprensione, uno scarto che le visualizzazioni statiche convenzionali, per i non specialisti, tendono ad allargare anziché a colmare.

Questa tesi si chiede in che misura le visualizzazioni esperienziali guidate dal machine learning migliorino la comprensione cognitiva dei sistemi climatici complessi a rischio di tipping tra i non specialisti, rispetto a quelle statiche convenzionali. Usa la Atlantic Meridional Overturning Circulation (AMOC) come caso di studio per un framework trasferibile costruito in tre fasi: \emph{rilevare} (\emph{detect}), un'analisi della struttura del sistema tramite machine learning; \emph{tradurre} (\emph{translate}), un artefatto interattivo che rende quella struttura mantenendo le sue condizioni statica ed esperienziale allineate per provenienza, inquadramento e unità di misura; e \emph{valutare} (\emph{evaluate}), uno studio esplorativo tra soggetti con venticinque non specialisti, le cui risposte scritte sono state analizzate rispetto a una griglia di codifica definita a priori.

I risultati sono informativi proprio perché non sono uniformi. Si è riscontrato che la presentazione esperienziale non sempre aumentava la comprensione. Ha aiutato in modo più netto dove la struttura di un sistema è temporale e invisibile in una figura statica, ha mostrato poco o nessun vantaggio dove una figura statica già presentava la struttura con chiarezza e, per una delle scene, ha reso una lettura errata sistematica più probabile anziché meno. I partecipanti stessi hanno nominato il meccanismo da entrambi i lati. Questi risultati sono stati ottenuti sull'AMOC e sui grafici più comunemente usati per analizzarla e monitorarla, perciò l'esito non è proposto come un risultato da trasferire altrove: ciò che la tesi avanza è il framework, non il risultato.

Il contributo è quindi principalmente metodologico: un framework, e un primo corpo di evidenze, per rendere intelligibili i sistemi invisibili --- non dimostrando che l'esperienza aiuti ovunque, ma mostrando per quali strutture, per quali osservatori e a quale costo percettivo.

\cleardoublepage
